\section{Balancing the torsion profiles}
\label{sec:construction}

We now construct $X$ so that the four fixed components of the
involution on $X\times Y$ have exactly the Morse count needed in
Theorem~\ref{thm:main}. The construction is explicit, and all homology
in this section has integral coefficients.

\subsection{Seven disjoint centers}
Set
\[
 N=1664,\qquad L=N-3=1661.
\]
The centers are chosen so that, in each of the degree lists $1+2j$
and $2+2j$, the difference between the $2$- and $3$-primary torsion
multiplicities is $-1$ for $1\leq j\leq L$, except for $+1$ at
$j=2,L-1$. Each positive entry is canceled by either adjacent
negative entry under convolution with $(1,1)$, whereas the middle
weight $6$ in $(1,6,1)$ leaves a positive excess. We first realize
this pattern and then compute the two counts.

Inside $\CP^N$, choose the following centers, with $S_2,S_3$ as in
Lemma~\ref{geo:surfaces}:
\begin{center}
\begin{tabular}{@{}lrrrr@{}}
\toprule
Center & Number & Vector dimension & Codimension & Total Betti rank\\
\midrule
$S_2\times\CP^1$ & $2$ & $90$ & $N-3$ & $16$\\
$S_3$ & $3$ & $165$ & $N-2$ & $8$\\
$S_3\times\CP^2$ & $2$ & $495$ & $N-4$ & $24$\\
\bottomrule
\end{tabular}
\end{center}
Here the vector dimension is one more than the dimension of the
receiving projective space. The embeddings of
Lemma~\ref{geo:surfaces}, followed by the Segre embeddings in the
product cases, give these dimensions. Since
\[
 2\cdot90+3\cdot165+2\cdot495=1665=N+1,
\]
we can place the seven centers in the projectivizations of distinct
summands of a direct-sum decomposition of $\C^{N+1}$. These projective
linear subspaces, and hence the centers, are pairwise disjoint.

Let $X$ be the successive complex blow-up along the seven centers.
The centers not yet blown up have unchanged neighborhoods because
they are disjoint. Their complex codimensions are at least two, so
Proposition~\ref{geo:blowup} shows that $X$ is simply connected,
closed, and K\"ahler. Its complex dimension is $N$.
For a center $S_p\times\CP^r$, the extra summands in $H_i(X)$ are
\begin{equation}
 \bigoplus_{k=1}^{L-r}H_{i-2k}(S_p\times\CP^r).
 \label{eq:center-shifts}
\end{equation}
Summing their free ranks gives
\begin{align}
 b(X)&=(N+1)+2\cdot16(L-1)+3\cdot8L+2\cdot24(L-2)
 \nonumber\\
 &=174281.
 \label{eq:rank-X-new}
\end{align}

\subsection{The two prime multiplicities}
The torsion of $S_p$ occurs in degrees $1$ and $2$, with the same
group in each degree. Accordingly, list the torsion degrees of $X$
as $1+2j$ and $2+2j$, with $j\in\Z$; all multiplicities are zero
outside their stated ranges. For $r=0,1,2$ define
\begin{equation}
 f_r(j)=\#\{(k,l)\in\Z^2:1\leq k\leq L-r,
                         \ 0\leq l\leq r,\ k+l=j\}.
 \label{eq:f-r}
\end{equation}
K\"unneth has no additional torsion terms here, since
$H_*(\CP^r;\Z)$ is free. Formula~\eqref{eq:center-shifts} therefore
gives, in either degree list,
\begin{equation}
 \Tor H_{1+2j}(X)\cong\Tor H_{2+2j}(X)
 \cong(\Z/2)^{a_j}\oplus(\Z/3)^{q_j},\qquad
 a_j=4f_1(j),\quad q_j=3f_0(j)+2f_2(j).
 \label{eq:new-profiles}
\end{equation}
The factor $4$ counts two centers, each contributing two copies
of $\Z/2$. Both sequences are supported on $1\leq j\leq L$.
Their values are
\begin{center}
\begin{tabular}{@{}lccc@{}}
\toprule
Indices & $1,L$ & $2,L-1$ & $3,\ldots,L-2$\\
\midrule
$f_0(j)$ & $1$ & $1$ & $1$\\
$f_1(j)$ & $1$ & $2$ & $2$\\
$f_2(j)$ & $1$ & $2$ & $3$\\
$a_j$ & $4$ & $8$ & $8$\\
$q_j$ & $5$ & $7$ & $9$\\
$\delta_j=a_j-q_j$ & $-1$ & $1$ & $-1$\\
\bottomrule
\end{tabular}
\end{center}
Thus $\delta$ is $-1$ throughout its support except at $2,L-1$,
where it is $1$. Writing $A=\sum_j a_j$ and $Q=\sum_jq_j$, we obtain
\begin{equation}
 A=8(L-1)=13280,\qquad Q=3L+6(L-2)=14937,
 \qquad A-Q=4-L<0.
 \label{eq:totals}
\end{equation}

\subsection{Products and the exact gap}
Choose any K\"ahler form $\omega_X$ on $X$. With the surface and
form of Proposition~\ref{inv:surface}, set
\[
 (M,\omega)=(X\times Y,\omega_X\oplus\omega_Y),
 \qquad C=X\times\CP^1.
\]
Both manifolds are simply connected, closed, and K\"ahler. Their
real dimensions are $3332$ and $3330$, respectively. The involution
$g=\id_X\times g_Y$ is Hamiltonian and has four fixed components
diffeomorphic to $C$.

For a finitely supported sequence $v=(v_j)_{j\in\Z}$ and a finite
sequence $e=(e_0,\ldots,e_s)$, write
\[
 (e*v)_j=\sum_{t=0}^s e_t v_{j-t}.
\]
The homology of the factors $\CP^1$ and $Y$ is free and concentrated
in even degrees. K\"unneth shows that the $2$- and $3$-primary
torsion multiplicities in either degree list of $C$ and $M$ are
$e*a$ and $e*q$, respectively, with
\begin{equation}
 e_C=(1,1),\qquad e_M=(1,6,1).
 \label{eq:kernels}
\end{equation}
The free ranks are $b(C)=2b(X)$ and $b(M)=8b(X)$.

For nonnegative integers $v,w$ one has
\begin{equation}
 d\bigl((\Z/2)^v\oplus(\Z/3)^w\bigr)=\max(v,w).
 \label{eq:prime-generators}
\end{equation}
Reduction modulo $2$ and modulo $3$ gives the lower bound; pairing
$\min(v,w)$ factors by the Chinese remainder theorem gives the
upper bound. After summing all degrees of either parity in $M$, the
prime multiplicities are $8A$ and $8Q$. Since $Q>A$,
\begin{equation}
 \beta_{\Z}^{\mathrm{cyc}}(M)=8b(X)+32Q.
 \label{eq:beta-value}
\end{equation}

The degreewise calculation is different. For $C$,
\[
 (e_C*\delta)_j=\delta_j+\delta_{j-1}\leq0
 \quad(j\in\Z),
\]
because each positive entry of $\delta$ has a negative entry on
either side. This also checks the boundary indices, where $\delta$
is zero outside $[1,L]$. Thus the $3$-primary multiplicity attains
the maximum in \eqref{eq:prime-generators} in every degree of $C$.
Its sum in each parity is $2Q$, and hence
\begin{equation}
 \lambda(C)=2b(X)+8Q,
 \qquad 4\lambda(C)=\beta_{\Z}^{\mathrm{cyc}}(M).
 \label{eq:fixedcomponent-count}
\end{equation}

For $M$, the positive entries of $e_M*\delta$ occur exactly at
$j=3,L$, and each has value $4$. Indeed,
\[
 (e_M*\delta)_3=\delta_3+6\delta_2+\delta_1=-1+6-1=4.
\]
Reflection gives the same positive value at $L$. At every other
index with $\delta_{j-1}=-1$, the three terms sum to at most
$1-6+1=-4$. If $\delta_{j-1}=0$, only nonpositive boundary entries
can remain. This proves that there are no other positive entries.
Figure~\ref{fig:profiles} displays the local
calculation at the left endpoint.

\begin{figure}[tb]
\centering
\begin{tikzpicture}[x=0.61cm,y=0.23cm,>=stealth]
\foreach \origin/\lab in {0/{\delta_j},5.6/{(e_C*\delta)_j},11.2/{(e_M*\delta)_j}} {
 \begin{scope}[xshift=\origin cm]
 \draw[->,gray] (-.5,0)--(5.7,0);
 \foreach \x in {0,...,5} {\node[below,font=\scriptsize] at (\x,-9.5) {\x};}
 \node at (2.5,6.8) {$\lab$};
 \end{scope}
}
\foreach \x/\v in {0/0,1/-1,2/1,3/-1,4/-1,5/-1} {
 \draw[very thick,blue!65!black] (\x,0)--(\x,\v);
 \fill[blue!65!black] (\x,\v) circle (2pt);
}
\node[above,font=\scriptsize] at (2,1.3) {$1$};
\node[below,font=\scriptsize] at (1,-1.3) {$-1$};
\begin{scope}[xshift=5.6cm]
\foreach \x/\v in {0/0,1/-1,2/0,3/0,4/-2,5/-2} {
 \draw[very thick,blue!65!black] (\x,0)--(\x,\v);
 \fill[blue!65!black] (\x,\v) circle (2pt);
}
\node[above,font=\scriptsize] at (2,.3) {$0$};
\node[below,font=\scriptsize] at (4,-2.3) {$-2$};
\end{scope}
\begin{scope}[xshift=11.2cm]
\foreach \x/\v in {0/0,1/-1,2/-5,3/4,4/-6,5/-8} {
 \draw[very thick,blue!65!black] (\x,0)--(\x,\v);
 \fill[blue!65!black] (\x,\v) circle (2pt);
}
\node[right,font=\small] at (3,4) {$4$};
\end{scope}
\end{tikzpicture}
\caption{The torsion difference near the left endpoint. Convolution
with $(1,1)$ removes its positive entry, whereas convolution with
$(1,6,1)$ produces an excess of $4$ at index $3$. The right endpoint
has the reflected pattern. Horizontal coordinates are the list
indices $j$, not homological degrees; all panels use the same
vertical scale.}
\label{fig:profiles}
\end{figure}

Using $x_+=\max(x,0)$, the degreewise generator sum in each parity
of $M$ is therefore
\[
 \sum_j\max\{(e_M*a)_j,(e_M*q)_j\}
 =8Q+\sum_j((e_M*\delta)_j)_+=8Q+8.
\]
There are two parity lists and a factor of two in $\lambda$, giving
\begin{equation}
 \lambda(M)=8b(X)+32Q+32
           =\beta_{\Z}^{\mathrm{cyc}}(M)+32.
 \label{eq:exact-gap}
\end{equation}

\subsection{Completion of the construction}
\begin{proof}[Proof of Theorem~\ref{thm:main}]
Apply Theorem~\ref{morse:smale} to $M$ and $C$, whose simple
connectivity and dimensions were verified above. It gives
$\SM(M)=\lambda(M)$ and an ordinary Morse function
$h:C\to\R$ with precisely $\lambda(C)$ critical points.
Corollary~\ref{inv:realization} gives a smooth one-periodic
Hamiltonian on $M$ with exactly $4\lambda(C)$ fixed points, all
nondegenerate and with contractible fixed-point loops. It also gives
$c_1(TM)(\pi_2(M))=\Z$.

In the explicit formula of Proposition~\ref{inv:perturb}, the
first-stage Hamiltonian can be taken autonomous: it is the pullback
of $k/2$, where $k$ generates the period-one circle action on $Y$.

Finally, \eqref{eq:rank-X-new}, \eqref{eq:totals},
\eqref{eq:fixedcomponent-count}, and \eqref{eq:exact-gap} give
\[
 \lambda(C)=468058,\qquad
 \beta_{\Z}^{\mathrm{cyc}}(M)=1872232,\qquad
 \SM(M)=1872264.
\]
Since $\dim_{\C}X=1664$ and $\dim_{\C}Y=2$, the dimension is
also as asserted.
\end{proof}
